% Tutorial set: 10; actual session: 2026-10-28 (Week 11).
% Sources: tute10.pdf p. 1, Q1; t10q1.ipynb, all 34 cells (zero-based 0--33).
% Adapted from templates/tutorial.tex. Code and cached outputs inspected;
% no model execution or training. Original notebook images are not reproduced.

\begingroup
\renewcommand{\recordid}[1]{}
\renewcommand{\relatedtopic}[2]{\hyperref[topic:#1]{#2}}
\exercisegroup{SC4001-TUT10}{Tutorial 10：Vision Transformer 推理流程}
\label{tutorial:SC4001-TUT10}
\noteintro{
\textbf{日期：}2026-10-28（第 11 周）。
\textbf{来源：}\texttt{tute10.pdf} p.~1，Q1；\texttt{t10q1.ipynb}，
署名 Hiroto Honda。已核对代码与文件保存的图表、输出；未重新运行模型。
本题使用预训练 ViT，围绕图像分类走读数据流，不涉及重新训练。
}

\exercisequestion{SC4001-TUT10-Q01}{从图像块到分类结果}
\nopagebreak[4]
\textbf{题目。}使用提供的 notebook 理解 Vision Transformer（ViT，视觉 Transformer）
的推理过程：图像切块与嵌入、位置编码、Transformer encoder（编码器）、
attention（注意力）可视化及分类头。

\textbf{对应讲义：}\relatedtopic{SC4001-L10-T02}{ViT 图像分类}、
\relatedtopic{SC4001-L09-T02}{自注意力与加权汇总}。

\textbf{1. 输入与图像块序列。}
\texttt{create\_model("vit\_base\_patch16\_224", pretrained=True)}
载入已有权重，对教堂示例图片预测 ImageNet 的 1000 类之一。
图片先调整至 $224\times224$，转为 RGB 通道张量；像素缩放至 $[0,1]$ 后，
按本例设置计算 $(x-0.5)/0.5$。\texttt{unsqueeze(0)} 增加 batch（批次）维。

每块大小为 $16\times16$，故块数为 $(224/16)^2=196$。
\texttt{model.patch\_embed} 使用如下卷积：
\begin{center}
\small\texttt{Conv2d(3, 768, kernel\_size=16, stride=16)}
\end{center}
每次覆盖一个不重叠的 patch（图像块），共享可学习投影产生 768 维表示，
再将空间位置排成序列。一个原始块含 $16\times16\times3=768$ 个数，
与本例 embedding（嵌入）维数相同，但投影仍会改变特征；一般模型不要求两者相等。

\begin{center}
\small
\begin{tabularx}{\textwidth}{@{}l l X@{}}
\toprule
阶段 & 形状 & 解释 \\
\midrule
输入 & $1\times3\times224\times224$ & 一张 RGB 图片 \\
切块投影 & $1\times768\times14\times14$ & 每个空间位置有 768 个特征 \\
排成序列 & $1\times196\times768$ & 196 个 patch tokens（图像块词元） \\
补齐输入 & $1\times197\times768$ & 前置 class token，再加位置编码 \\
12 个编码块 & $1\times197\times768$ & 保留位置数与特征维，逐层更新内容 \\
读取 class token & $1\times768$ & 取最终序列的第 0 行 \\
分类头 & $1\times1000$ & 每个类别一个 logit（未归一化分数） \\
\bottomrule
\end{tabularx}
\end{center}

\textbf{2. 位置与整图表示。}
输入由以下代码构成（本例 batch 大小为 1）：
\begin{center}
\small\ttfamily
torch.cat((model.cls\_token, patches), dim=1) + pos\_embed
\end{center}
\texttt{cat} 将可学习的 class token（分类词元）放到序列最前面，
长度从 196 增至 197；位置嵌入逐元素相加，形状保持不变。
忽略 batch 维，写成
\[
 Z_0=[c;\,p_1;\,\ldots;\,p_{196}]+E_{\rm pos},
 \qquad Z_0,E_{\rm pos}\in\mathbb R^{197\times768}.
\]
$p_i$ 表示第 $i$ 块内容，位置向量提供其空间身份。初始 $c$ 是各图像共享的参数，
不含已知类别标签；经自注意力更新后，它汇总当前图像的信息，供分类头读取。

\textbf{位置编码图的读法。}代码固定一个 patch 的位置向量，计算它与
全部 196 个 patch 位置向量的 cosine similarity（余弦相似度），
再排回 $14\times14$。这些小图展示学到的位置关系，不依赖当前示例图像的内容；
邻近位置可能呈现较高相似度，但这不是训练机制强制保证的性质。

\textbf{3. 编码块与多头自注意力。}
本模型有 12 个串联的 encoder blocks，每个 block 内有 12 个并行的 heads（注意力头）。
所有头均读取全部 197 个 token；每个头的特征维数为 $768/12=64$。
按 notebook 引用的实现，块内采用 pre-norm（前归一化）：
\[
 U=Z+\operatorname{MSA}(\operatorname{LN}(Z)),\qquad
 Z_{\rm next}=U+\operatorname{MLP}(\operatorname{LN}(U)).
\]
LN 为层归一化；MSA 在 token 之间交换信息；MLP（多层感知机）对每个 token
使用同一组参数加工特征；残差相加保留原表示并加入更新。
全图已知，分类使用 encoder，无需 decoder 或 causal mask（因果掩码）。

QKV 投影的输出维为 $3\times768=2304$。省略 batch 维，代码先得到
$197\times2304$，再重排为 $197\times3\times12\times64$，
分出 Q、K、V 并调整轴顺序。对头 $h$：
\[
 Q_h,K_h,V_h\in\mathbb R^{197\times64},\quad
 A_h=\operatorname{softmax}_{\rm row}\!\left(\frac{Q_hK_h^\top}{8}\right),
 \quad O_h=A_hV_h\in\mathbb R^{197\times64}.
\]
$A_h$ 的行对应 query（查询），列对应 key（键）；每行权重对所有 value（值）
加权汇总，生成该位置的新表示。12 个头的输出拼接为 $197\times768$，
经输出投影后进入残差路径。维度划分发生在特征轴，不是把 token 分给不同的头。

\textbf{4. 从完整分数矩阵读出空间图。}
Notebook 的 \texttt{attention\_matrix = q @ kT} 得到
$12\times197\times197$，其中第一个轴为 head。
\texttt{attention\_matrix[3]} 显示第 4 个头的完整匹配矩阵。
空间热图则读取 \texttt{attention\_matrix[i,100,1:]}：
固定查询 token 的索引为 100，去掉第 0 列的 class token，将剩下的 196 个分数排回图像位置。

\begin{center}
% Original indexing diagram based on t10q1.ipynb, cells 29--30 (zero-based).
% Highlighting indicates array selection only; no model scores are plotted.
\begin{tikzpicture}[font=\small,>=Stealth]
  \fill[noteBlue!10] (0,0) rectangle (0.24,2.1);
  \fill[noteBlue!35] (0,1) rectangle (2.6,1.18);
  \draw (0,0) rectangle (2.6,2.1);
  \draw[noteBlue] (0.24,0) -- (0.24,2.1);
  \node[align=center] at (1.3,2.55) {固定头 $i$\\$197\times197$};
  \node[anchor=east,font=\scriptsize] at (0,1.09) {行 100};
  \node[anchor=north,font=\scriptsize,align=center] at (1.3,-0.12)
    {列 0：class token\\列 1--196：patches};
  \draw[->,thick,noteBlue] (2.8,1.09) -- (4.15,1.09)
    node[midway,above,font=\scriptsize] {取 \texttt{1:}};
  \filldraw[fill=noteBlue!25,draw=noteBlue] (4.3,0.99) rectangle (6.9,1.19);
  \node[align=center] at (5.6,1.72) {196 个分数};
  \node[align=center,font=\scriptsize] at (5.6,0.45) {保持原来的 patch 顺序};
  \draw[->,thick,noteBlue] (7.05,1.09) -- (8.45,1.09)
    node[midway,above,font=\scriptsize] {reshape};
  \begin{scope}[shift={(8.65,0.04)}]
    \fill[noteBlue!6] (0,0) rectangle (2.1,2.1);
    \draw[step=0.15,gray!60,thin] (0,0) grid (2.1,2.1);
    \draw (0,0) rectangle (2.1,2.1);
  \end{scope}
  \node[align=center] at (9.7,2.55) {$14\times14$\\patch 位置图};
  \node[align=center,font=\scriptsize] at (9.7,-0.35) {每格对应一个 key patch};
\end{tikzpicture}

\end{center}
{\small 原创索引示意，依据 notebook 的可视化代码；图中颜色仅标识选取位置，不表示分数大小。}

由于序列第 0 项是 class token，查询索引 100 对应从零计数的 patch 99，
即 $14\times14$ 网格的第 8 行、第 2 列。
\texttt{range(7)} 展示头 0--6 对同一查询位置的匹配分数，不能直接当作最终分类的归因图。

\textbf{可视化勘误。}该手工分支既未先经过实际 block 的 \texttt{norm1}，
也未对 $QK^\top$ 进行缩放和 softmax。因此保存的图是该手工计算的原始匹配分数，
可能为负、每行和不必为 1；不能将其标作模型实际的归一化注意力权重。
上述正确计算顺序可查 notebook 引用的
\href{https://github.com/huggingface/pytorch-image-models/blob/198f6ea0f3dae13f041f3ea5880dd79089b60d61/timm/models/vision_transformer.py}
{\texttt{timm} 源码中的 \texttt{Block.forward} 与 \texttt{Attention.forward}}。

\textbf{5. 读取 class token 并分类。}
\texttt{for blk in model.blocks: x = blk(x)} 依次执行 12 个块；
随后 \texttt{x = model.norm(x)}，再以 \texttt{x[:,0]} 取出 class token 的最终表示。
此处选择一行，不是将 197 行展平或直接求和。
文件打印的 \texttt{model.head} 是 \texttt{Linear(768,1000)}，计算
\[
 s=h_{\rm CLS}W+b\in\mathbb R^{1\times1000},\qquad
 \widehat y=\arg\max_k s_k.
\]
\texttt{imagenet\_labels} 将类别索引映射为名称。
分数图的横轴为类别编号、纵轴为 logit；若需要概率，再对 1000 个分数做 softmax，
最大分数对应的类别保持不变。教程称其为 MLP head，但本例实际分类头只有一个线性层，
应与 encoder 内的 MLP 区分。

\textbf{保存结果的边界。}完整调用 \texttt{model(img\_tensor)} 的缓存标签为
\texttt{church, church\_building}；手工拆解分支的缓存标签为
\texttt{dome}（编号 538）。两者不一致，仅凭该文件无法判定原因；
若要验证拆解等价性，应在一致的模型、输入与推理状态下重跑比较。
本次仅核对已有文件，不把缓存输出当作重新复现成功的证据。
\endgroup
